% HLCC Loss on Small Networks
% Compile with: pdflatex small_network_paper.tex && bibtex small_network_paper && pdflatex small_network_paper && pdflatex small_network_paper

\documentclass[11pt,a4paper]{article}

% Packages
\usepackage[utf8]{inputenc}
\usepackage[T1]{fontenc}
\usepackage{amsmath,amssymb,amsthm}
\usepackage{booktabs}
\usepackage{graphicx}
\usepackage{hyperref}
\usepackage{cleveref}
\usepackage{algorithm}
\usepackage{algorithmic}
\usepackage{multirow}
\usepackage{subcaption}
\usepackage[margin=2.5cm]{geometry}
\usepackage{enumitem}
\usepackage{xcolor}
\usepackage{float}

% Theorem environments
\newtheorem{theorem}{Theorem}
\newtheorem{proposition}{Proposition}
\newtheorem{corollary}{Corollary}
\newtheorem{definition}{Definition}
\newtheorem{remark}{Remark}

% Custom commands (shared with hlcc_paper)
\newcommand{\hlcc}{\textsc{hlcc}}
\newcommand{\E}{\mathbb{E}}
\newcommand{\R}{\mathbb{R}}
\newcommand{\cstar}{c^{*}}
\DeclareMathOperator*{\argmax}{arg\,max}
\DeclareMathOperator*{\argmin}{arg\,min}

% Data repository URL — update when paper repo is made public
\newcommand{\datarepo}{https://github.com/USER/NNConfidence/blob/main}

\title{HLCC Loss on Small Networks:\\Analysing Normalisation-Aware Confidence\\in a Controlled Setting}
\author{
  Simon McCallum\\
  \texttt{[email]}
}
\date{\today}

\begin{document}
\maketitle

% ============================================================
\begin{abstract}
The Hybrid Linear-Convex Confidence (\hlcc{}) loss was originally proposed
for calibrating large language model (LLM) confidence heads, where the
opacity of billion-parameter networks makes it difficult to inspect
\emph{why} the method works.  In this paper, we isolate the core
mechanisms of \hlcc{} --- the asymmetric scoring function, the
normalisation-derived doubt signals, and the trainable confidence head
--- in a small, fully inspectable four-layer feedforward network
(20--24--16--5).  Using a synthetic concentric-rings classification task
with controlled label noise and out-of-distribution (OOD) regions, we
run a factorial experiment across 13 training configurations varying
the loss function (cross-entropy, \hlcc{}, combined, focal), normalisation
(present/absent), and confidence head state (frozen at 0.5 vs.\
trainable).  We analyse weight structure, activation distributions,
doubt signal informativeness, and calibration quality, providing
fine-grained evidence for which components of the \hlcc{} system
contribute to improved confidence estimation.  Critically, confidence
is derived from raw output logits rather than softmax probabilities,
testing whether the \hlcc{} framework can operate without the
probabilistic assumptions underlying standard calibration methods.
\end{abstract}

% ============================================================
\section{Introduction}
\label{sec:intro}

Confidence calibration --- ensuring that a model's stated certainty
matches its empirical accuracy --- is a long-standing problem in
machine learning~\cite{guo2017calibration, naeini2015obtaining}.
Recent work has extended calibration to LLMs through methods such as
maximum softmax probability, verbalized confidence, semantic entropy,
and self-consistency~\cite{kadavath2022language, xiong2024llms,
kuhn2023semantic, wang2023selfconsistency}.

The \hlcc{} loss~\cite{mccallum2026hlcc}, inspired by confidence-based
assessment in education~\cite{mccallum2010gamedesign,
gardnermedwin1995}, proposes an asymmetric scoring function where
correct predictions receive a linear reward $R(c) = 1 + c$ and
incorrect predictions incur a quadratic penalty $P(c) = -2c^2$.  When
combined with a \emph{NormShift} confidence head that routes
layer normalisation dynamics to a learned confidence estimator, this
system achieves strong discrimination (AUROC up to 0.87) and
calibration (ECE as low as 0.105) on LLMs with 7--8 billion
parameters.

However, at LLM scale, it is difficult to determine \emph{which
component} drives the improvement.  Is it the asymmetric loss
function?  The normalisation signals?  The confidence head
architecture?  The interaction between them?  To answer these
questions, we construct a minimal testbed: a four-layer feedforward
network small enough to inspect every weight, activation, and
gradient.

\paragraph{Research questions.}  We address four specific questions:
\begin{enumerate}[label=\textbf{Q\arabic*:}]
  \item Does \hlcc{} improve calibration over cross-entropy in small
        networks?
  \item Does layer normalisation provide useful doubt signals at
        small scale?
  \item Can a confidence head learn from raw output logits (without
        softmax)?
  \item How do network weights differ between normalised and
        non-normalised variants trained with the same loss?
\end{enumerate}

% ============================================================
\section{Background}
\label{sec:background}

\subsection{HLCC Scoring Function}

The \hlcc{} scoring function assigns a score to each prediction based
on the model's confidence $c \in [0, 1]$ and the correctness
indicator $y \in \{0, 1\}$:
\begin{equation}
  S(c, y) = \begin{cases}
    1 + c & \text{if } y = 1 \quad \text{(correct)} \\
    -2c^2 & \text{if } y = 0 \quad \text{(incorrect)}
  \end{cases}
  \label{eq:hlcc}
\end{equation}

The optimal confidence under this scoring rule for a question with
true accuracy $p$ is:
\begin{equation}
  \cstar = \frac{p}{4(1-p)}
  \label{eq:optimal}
\end{equation}
which naturally suppresses overconfidence: even at $p = 0.8$, the
optimal confidence is only $\cstar = 1.0$~\cite{mccallum2026hlcc}.

\subsection{Normalisation as Uncertainty Signal}

Layer normalisation~\cite{ba2016layer} rescales activations to zero
mean and unit variance.  The \emph{magnitude} of this correction
carries information: when pre-normalisation activations have unusual
scale or variance, the normalisation layer must apply large
corrections, which may indicate the model is ``forcing'' a
commitment~\cite{mccallum2026hlcc}.

The \texttt{UncertaintyAwareLayerNorm} module captures three
\emph{doubt signals} per layer:
\begin{enumerate}
  \item \textbf{Pre-normalisation variance}: high variance may indicate
        confident separation between features.
  \item \textbf{Correction scale}: $(v + \epsilon)^{-1/2}$, measuring how
        much the normalisation rescaled the activations.
  \item \textbf{Pre-normalisation L2 norm}: the overall magnitude of
        activations before normalisation.
\end{enumerate}

\subsection{Confidence Without Softmax}

Standard calibration methods derive confidence from probability
distributions (softmax outputs, logit-based entropy, etc.).  Here
we deliberately avoid softmax and instead derive confidence from
raw output logits:
\begin{equation}
  c_{\text{raw}} = \sigma\!\left(\max_i z_i - \frac{1}{C-1}\sum_{j \neq i^*} z_j\right)
  \label{eq:raw_conf}
\end{equation}
where $z_i$ are the output logits, $i^* = \argmax_i z_i$, $C$ is the
number of classes, and $\sigma$ is the sigmoid function.  This tests
whether the \hlcc{} framework can operate on activation geometry
rather than probability estimates.

% ============================================================
\section{Experimental Design}
\label{sec:design}

\subsection{Network Architecture}

We use a four-layer feedforward network with the following structure:
\begin{equation}
  \mathbf{x} \in \R^{20} \xrightarrow{\text{FC}} \R^{24}
  \xrightarrow{[\text{LN}]} \xrightarrow{\text{ReLU}}
  \xrightarrow{\text{FC}} \R^{16}
  \xrightarrow{[\text{LN}]} \xrightarrow{\text{ReLU}}
  \xrightarrow{\text{FC}} \R^{5}
  \label{eq:arch}
\end{equation}
where $[\text{LN}]$ denotes optional \texttt{UncertaintyAwareLayerNorm}
layers.  The network has approximately 1,000 parameters, making it
possible to inspect every weight.

The confidence head takes the second hidden layer activations
($\R^{16}$) concatenated with 6 doubt signals (3 per normalisation
layer) and maps them through a single hidden layer of size 8 to a
scalar confidence in $[0, 1]$:
\begin{equation}
  c = \sigma\!\left(\mathbf{w}_2^{\top} \text{ReLU}(\mathbf{W}_1 [\mathbf{h}_2; \mathbf{d}] + \mathbf{b}_1) + b_2\right)
  \label{eq:head}
\end{equation}
The head is initialised to output $c \approx 0.5$ (via zero biases and
small weights), ensuring that the frozen-confidence variants start at
the neutral point of the \hlcc{} scoring function.

\subsection{Datasets}
\label{sec:datasets}

We evaluate on six datasets spanning synthetic and real-world tasks,
varying in difficulty, dimensionality, and number of classes.

\paragraph{Concentric Rings (primary synthetic).}
A 5-class classification task using concentric annular rings in 2D,
embedded in 20 dimensions (18 uninformative noise features scaled by
0.3).  We run two noise settings: the \emph{standard} variant (15\%
label noise, $\sim$93\% peak accuracy) and a \emph{hard} variant
(30\% noise, $\sim$75\% accuracy).  The hard variant places accuracy
in the 50--80\% range where the \hlcc{} scoring function's
conservatism differs meaningfully from standard calibration methods
--- at 93\% accuracy, the \hlcc{} optimal confidence saturates at
$\cstar = 1.0$, eliminating the asymmetry that distinguishes \hlcc{}
from Brier scoring.

\paragraph{Two Moons (standard benchmark).}
Two interleaving half-circles (\texttt{make\_moons}, noise$=0.2$),
the de facto standard for visualising uncertainty
surfaces~\cite{liu2020simple}.  Binary classification in 2D with
3{,}000 training points.  OOD data comprises shifted moons
(translated by $(3,3)$) and high-noise moons (noise$=0.8$).

\paragraph{UCI Breast Cancer Wisconsin.}
A real-world binary classification task (569 samples, 30 features)
from the UCI repository~\cite{dua2019uci}.  Standardised on training
set statistics.  OOD data is created by Gaussian perturbation of test
features ($\sigma = 2\times$ feature std).

\paragraph{Sklearn Digits.}
8$\times$8 pixel digit images (1{,}797 samples, 10 classes).
Tests scaling to more classes with a small network.

\paragraph{Overlapping Blobs.}
Synthetic multi-class task via \texttt{make\_classification} with
$\texttt{class\_sep}=0.8$, 8 informative features, 4 redundant, and
10\% label noise on training data only.  5 classes in 20 dimensions,
matching the concentric rings dimensionality for controlled comparison.

All datasets use the same split protocol: separate train/validation/test
with clean labels on validation and test sets.  OOD sets use either
geometric construction (rings), distributional shift (moons, blobs),
or Gaussian perturbation (UCI, digits).

\subsection{Calibration Baselines}
\label{sec:baselines}

We compare the \hlcc{} confidence head against four standard
calibration methods:

\begin{enumerate}
  \item \textbf{Temperature scaling}~\cite{guo2017calibration}: A
        single scalar $T$ is fit on the validation set to minimise
        NLL of $\text{softmax}(\mathbf{z}/T)$.  Applied post-hoc to
        every trained model.
  \item \textbf{Focal loss}~\cite{lin2017focal}: An alternative
        asymmetric loss $\text{FL}(p_t) = -(1-p_t)^\gamma \log p_t$
        ($\gamma = 2$) that down-weights easy examples.  Included to
        answer the natural question: ``Is \hlcc{} just a fancy focal
        loss?''
  \item \textbf{MC Dropout}~\cite{gal2016dropout}: Dropout ($p=0.1$)
        applied to hidden activations at test time across 30 forward
        passes.  Predictive entropy provides the uncertainty estimate.
  \item \textbf{Deep Ensemble}~\cite{lakshminarayanan2017simple}: 5
        independently initialised networks trained with CE loss.
        Averaged softmax probabilities provide calibrated confidence;
        disagreement measures uncertainty.  The gold standard for
        small-network uncertainty.
\end{enumerate}

\subsection{Statistical Protocol}
\label{sec:stats}

All experiments are run across 5 random seeds (varying model
initialisation, fixed dataset).  We report mean $\pm$ standard
deviation for all metrics.  This addresses a common weakness in
calibration papers where single-run results may reflect
initialisation luck rather than method quality.

We report three ECE variants: standard equal-width ECE (15 bins),
debiased ECE~\cite{kumar2019verified} (which corrects for finite-sample
bias), and adaptive ECE (equal-mass bins, robust to skewed confidence
distributions).

\subsection{Training Variants}
\label{sec:variants}

We train 13 configurations spanning all combinations of three factors
plus focal loss baselines (\Cref{tab:variants}).

\begin{table}[H]
\centering
\caption{Training variant matrix.  Each row is a distinct configuration.
Variants 11--13 use focal loss ($\gamma=2$) as the base classification loss.}
\label{tab:variants}
\begin{tabular}{@{}rlccl@{}}
\toprule
\# & Name & Loss & Norm & Confidence \\
\midrule
1  & \texttt{ce\_nonorm}          & CE        & No  & Fixed 0.5 \\
2  & \texttt{ce\_norm}            & CE        & Yes & Fixed 0.5 \\
3  & \texttt{hlcc\_nonorm\_frozen}    & HLCC      & No  & Fixed 0.5 \\
4  & \texttt{hlcc\_norm\_frozen}      & HLCC      & Yes & Fixed 0.5 \\
5  & \texttt{hlcc\_nonorm\_trainable} & HLCC      & No  & Trainable \\
6  & \texttt{hlcc\_norm\_trainable}   & HLCC      & Yes & Trainable \\
7  & \texttt{combined\_nonorm}    & CE+HLCC   & No  & Trainable \\
8  & \texttt{combined\_norm}      & CE+HLCC   & Yes & Trainable \\
9  & \texttt{ce\_norm\_trainable}     & CE        & Yes & Trainable \\
10 & \texttt{hlcc\_norm\_2phase}      & CE$\to$HLCC & Yes & Frozen$\to$Trainable \\
\midrule
11 & \texttt{focal\_nonorm}      & Focal     & No  & Fixed 0.5 \\
12 & \texttt{focal\_norm}        & Focal     & Yes & Fixed 0.5 \\
13 & \texttt{focal\_hlcc\_norm}  & Focal+HLCC & Yes & Trainable \\
\bottomrule
\end{tabular}
\end{table}

Variant 10 uses a two-phase schedule: 50 epochs of pure
cross-entropy with frozen confidence (phase 1), then pure \hlcc{}
with trainable confidence (phase 2).

All variants are trained with Adam ($\eta = 10^{-3}$) for up to 200
epochs with early stopping (patience 20) on validation loss.

\subsection{Confidence Estimation}

For variants with a frozen confidence head, $c = 0.5$ for all
examples, representing maximum uncertainty under the \hlcc{} scoring
function (since $S(0.5, 1) = 1.5$ and $S(0.5, 0) = -0.5$).

For variants with a trainable confidence head, $c$ is learned
end-to-end via backpropagation through the \hlcc{} loss.  The
confidence head receives both semantic information (hidden layer
activations) and normalisation dynamics (doubt signals), when
normalisation is present.

We additionally compute raw-output confidence
(\Cref{eq:raw_conf}) as a non-learned baseline for all variants.

\subsection{Metrics}

We evaluate on the following metrics:
\begin{itemize}
  \item \textbf{ECE} (Expected Calibration Error)~\cite{naeini2015obtaining}:
        average gap between confidence and accuracy across 10 bins.
        Lower is better.
  \item \textbf{Brier Score}~\cite{guo2017calibration}:
        mean squared error between confidence and correctness.
        Lower is better.
  \item \textbf{HLCC Score}: mean $S(c, y)$ across all examples.
        Higher is better.
  \item \textbf{AUROC}: area under the ROC curve treating confidence as a
        discriminator for correct/incorrect predictions.  Higher is better.
  \item \textbf{AUPRC}: area under the precision-recall curve.
        Higher is better.
  \item \textbf{PRR} (Prediction Rejection Ratio): normalised area
        between the model's rejection curve and random
        rejection~\cite{fadeeva2024lmpolygraph}.  Higher is better.
  \item \textbf{Acc@$k$\%}: accuracy on the $k$\% most confident
        predictions (selective prediction).
\end{itemize}

% ============================================================
\section{Results}
\label{sec:results}

\subsection{In-Distribution Performance}

% TABLE 2: In-distribution results (rings dataset, mean over 5 seeds)
\begin{table}[H]
\centering
\caption{In-distribution test set results for all 13 variants on the rings dataset
(mean over 5 seeds).  Variants below the mid-rule use focal loss ($\gamma=2$) as the base classifier.}
\label{tab:indist}
\begin{tabular}{@{}lccccccc@{}}
\toprule
Variant & Acc & ECE$\downarrow$ & Brier$\downarrow$ & HLCC$\uparrow$ & AUROC$\uparrow$ & PRR$\uparrow$ & Conf \\
\midrule
\texttt{ce\_nonorm}                      & .937 & .437 & .250 & 1.37 & .505 & $-$.008 & .500 \\
\texttt{ce\_norm}                        & .920 & .420 & .250 & 1.34 & .511 & .031 & .500 \\
\texttt{hlcc\_nonorm\_frozen}            & .933 & .433 & .250 & 1.37 & .513 & .021 & .500 \\
\texttt{hlcc\_norm\_frozen}              & .915 & .415 & .250 & 1.33 & .504 & .012 & .500 \\
\texttt{hlcc\_nonorm\_trainable}         & .920 & .062 & .077 & 1.66 & .584 & .244 & .937 \\
\texttt{hlcc\_norm\_trainable}           & .892 & .086 & .099 & 1.58 & \textbf{.665} & \textbf{.382} & .930 \\
\texttt{combined\_nonorm}                & \textbf{.937} & \textbf{.034} & \textbf{.060} & \textbf{1.73} & .514 & .006 & .968 \\
\texttt{combined\_norm}                  & .918 & .076 & .081 & 1.68 & .600 & .205 & .970 \\
\texttt{ce\_norm\_trainable}             & .920 & .418 & .248 & 1.34 & .548 & .158 & .502 \\
\texttt{hlcc\_norm\_2phase}              & .908 & .077 & .085 & 1.63 & .639 & .294 & .949 \\
\midrule
\texttt{focal\_nonorm}                   & .921 & .421 & .250 & 1.34 & .488 & $-$.035 & .500 \\
\texttt{focal\_norm}                     & .917 & .417 & .250 & 1.33 & .516 & .051 & .500 \\
\texttt{focal\_hlcc\_norm}               & .910 & .054 & .083 & 1.61 & .603 & .189 & .904 \\
\bottomrule
\end{tabular}
\par\smallskip{\footnotesize\textit{Source:} \texttt{\datarepo/data/results/small\_network/results\_multi\_seed\_rings.json}}
\end{table}

\subsection{Confidence Calibration}

The frozen-confidence variants (1--4, 9) all produce $c = 0.5$ for
every example, yielding ECE $\approx$ accuracy $- 0.5$
regardless of loss or normalisation.  This confirms that calibration
requires a \emph{trainable} confidence signal.

Among trainable variants on the rings dataset, \texttt{combined\_nonorm}
achieves the lowest ECE (0.034) and Brier score (0.060), followed by
\texttt{focal\_hlcc\_norm} (ECE 0.054, Brier 0.083) and
\texttt{combined\_norm} (ECE 0.076, Brier 0.081).  This calibration
improvement is consistent across all six datasets
(\Cref{app:full}): on two\_moons, the best ECE reaches 0.020; on
breast\_cancer, 0.023; on digits, 0.042.  Even on the hardest task
(blobs, $\sim$72\% accuracy), \texttt{hlcc\_norm\_2phase} achieves
ECE 0.068 vs.\ 0.241 for the CE baseline.

The CE-trained variant with trainable confidence
(\texttt{ce\_norm\_trainable}) shows almost no calibration
improvement across any dataset (ECE 0.418--0.482), since the CE loss
does not provide gradients through the confidence head.

% Reliability diagrams will be referenced here
% \begin{figure}[H]
% \centering
% \includegraphics[width=\textwidth]{../data/results/small_network/plots/reliability_combined_norm.png}
% \caption{Reliability diagrams for selected variants.}
% \label{fig:reliability}
% \end{figure}

\subsection{Out-of-Distribution Detection}

% TABLE 3: OOD results (rings dataset, mean over 5 seeds)
\begin{table}[H]
\centering
\caption{OOD test set results on rings dataset (mean over 5 seeds).  Softmax AUROC
is the AUROC of the softmax-derived confidence (unsupervised baseline).}
\label{tab:ood}
\begin{tabular}{@{}lcccccc@{}}
\toprule
Variant & Acc & ECE$\downarrow$ & HLCC$\uparrow$ & AUROC$\uparrow$ & Conf & Softmax AUROC \\
\midrule
\texttt{ce\_nonorm}                      & .283 & .217 & 0.07 & .849 & .500 & .165 \\
\texttt{ce\_norm}                        & .288 & .212 & 0.08 & .850 & .500 & .165 \\
\texttt{hlcc\_nonorm\_frozen}            & .286 & .214 & 0.07 & .856 & .500 & .171 \\
\texttt{hlcc\_norm\_frozen}              & .292 & .208 & 0.08 & .848 & .500 & .166 \\
\texttt{hlcc\_nonorm\_trainable}         & .296 & .660 & $-$0.73 & .423 & .938 & .155 \\
\texttt{hlcc\_norm\_trainable}           & .281 & .681 & $-$0.81 & .218 & .927 & .165 \\
\texttt{combined\_nonorm}                & .286 & .677 & $-$0.78 & .626 & .963 & .175 \\
\texttt{combined\_norm}                  & .291 & .689 & $-$0.80 & .248 & .954 & .165 \\
\texttt{ce\_norm\_trainable}             & .288 & .223 & 0.07 & .409 & .502 & .165 \\
\texttt{hlcc\_norm\_2phase}              & .290 & .663 & $-$0.78 & .262 & .939 & .162 \\
\midrule
\texttt{focal\_nonorm}                   & .283 & .217 & 0.07 & .855 & .500 & .164 \\
\texttt{focal\_norm}                     & .290 & .210 & 0.08 & .849 & .500 & .160 \\
\texttt{focal\_hlcc\_norm}               & .289 & .622 & $-$0.66 & .442 & .892 & .161 \\
\bottomrule
\end{tabular}
\par\smallskip{\footnotesize\textit{Source:} \texttt{\datarepo/data/results/small\_network/results\_multi\_seed\_rings.json}}
\end{table}

\subsection{Effect of Normalisation}

On the rings dataset, normalisation consistently improves discrimination
(AUROC) when the confidence head is trainable:
\texttt{hlcc\_norm\_trainable} achieves AUROC 0.665 vs.\ 0.584 for
\texttt{hlcc\_nonorm\_trainable}, and \texttt{combined\_norm} achieves
0.600 vs.\ 0.514 for \texttt{combined\_nonorm}.  This pattern holds on
rings\_hard (0.662 vs.\ 0.588) and blobs (0.654 vs.\ 0.641).

However, on some datasets the normalisation effect reverses:
on breast\_cancer, \texttt{hlcc\_nonorm\_trainable} achieves AUROC
0.909 vs.\ 0.896 for the normalised variant, and on digits,
\texttt{hlcc\_nonorm\_trainable} achieves AUROC 0.756 vs.\ only 0.506
for \texttt{hlcc\_norm\_trainable}.  This suggests that the benefit of
normalisation-derived doubt signals depends on whether the task's
difficulty structure creates informative normalisation dynamics ---
when the network easily separates classes, normalisation statistics
carry less uncertainty information.

Normalised variants show slightly lower raw accuracy across all
datasets (e.g.\ 0.892--0.918 vs.\ 0.920--0.937 on rings), suggesting
a persistent trade-off between accuracy and confidence quality.

The PRR metric shows the clearest normalisation benefit on
rings-family datasets: \texttt{hlcc\_norm\_trainable} achieves PRR
0.382 vs.\ 0.244 on rings, and 0.362 vs.\ 0.230 on rings\_hard.

\subsection{Selective Prediction}

% TABLE 4: Selective prediction (rings dataset, mean over 5 seeds)
\begin{table}[H]
\centering
\caption{Selective prediction on rings dataset (mean over 5 seeds): accuracy when
retaining the top $k$\% of examples ranked by confidence.}
\label{tab:selective}
\begin{tabular}{@{}lccccc@{}}
\toprule
Variant & Acc@70\% & Acc@80\% & Acc@90\% & PRR$\uparrow$ & AUPRC$\uparrow$ \\
\midrule
\texttt{ce\_nonorm}                      & .939 & .940 & .936 & $-$.008 & .936 \\
\texttt{ce\_norm}                        & .923 & .925 & .921 & .031 & .923 \\
\texttt{hlcc\_nonorm\_frozen}            & .936 & .936 & .933 & .021 & .934 \\
\texttt{hlcc\_norm\_frozen}              & .919 & .920 & .917 & .012 & .917 \\
\texttt{hlcc\_nonorm\_trainable}         & .931 & .934 & .927 & .244 & \textbf{.945} \\
\texttt{hlcc\_norm\_trainable}           & .931 & .924 & .914 & \textbf{.382} & .934 \\
\texttt{combined\_nonorm}                & .940 & .941 & .939 & .006 & .938 \\
\texttt{combined\_norm}                  & .935 & .942 & .930 & .205 & .933 \\
\texttt{ce\_norm\_trainable}             & .929 & .927 & .924 & .158 & .935 \\
\texttt{hlcc\_norm\_2phase}              & .935 & .941 & .926 & .294 & .934 \\
\midrule
\texttt{focal\_nonorm}                   & .922 & .923 & .919 & $-$.035 & .918 \\
\texttt{focal\_norm}                     & .920 & .920 & .917 & .051 & .921 \\
\texttt{focal\_hlcc\_norm}               & .929 & .932 & .926 & .189 & .926 \\
\bottomrule
\end{tabular}
\par\smallskip{\footnotesize\textit{Source:} \texttt{\datarepo/data/results/small\_network/results\_multi\_seed\_rings.json}}
\end{table}

% ============================================================
\section{Analysis}
\label{sec:analysis}

\subsection{Weight Structure Comparison}

% TABLE 5: Weight analysis
\begin{table}[H]
\centering
\caption{Weight matrix statistics for normalised vs.\ non-normalised variants
(both trained with combined CE+HLCC loss).}
\label{tab:weights}
\begin{tabular}{@{}lcccccc@{}}
\toprule
Parameter & \multicolumn{3}{c}{With Norm} & \multicolumn{3}{c}{No Norm} \\
\cmidrule(lr){2-4} \cmidrule(lr){5-7}
 & $|\bar{w}|$ & Spectral & Eff.\ Rank & $|\bar{w}|$ & Spectral & Eff.\ Rank \\
\midrule
\texttt{layer1.weight} & 0.082 & 0.98 & 15.7 & 0.149 & 1.86 & 16.0 \\
\texttt{layer2.weight} & 0.107 & 1.73 & 12.2 & 0.277 & 6.04 & 9.5 \\
\texttt{output.weight} & 0.337 & 2.69 & 4.3  & 0.468 & 4.70 & 3.9 \\
\bottomrule
\end{tabular}
\end{table}

\subsection{Doubt Signals as Uncertainty Indicators}

When \texttt{UncertaintyAwareLayerNorm} is present, each normalisation
layer produces three doubt signals per example.  Comparing these
signals between correctly and incorrectly classified examples reveals
that pre-normalisation variance and L2 norm show statistically
significant differences between the two groups, confirming that
normalisation dynamics carry uncertainty information even at small
scale.  The second layer's signals are more informative than the
first layer's, consistent with the hypothesis that later layers
capture higher-level uncertainty about the classification decision.

% Doubt signal violin plots referenced here

\subsection{Discrimination: Trained Head vs.\ Unsupervised Baselines}
\label{sec:discrimination}

We compare the trained confidence head's AUROC against the unsupervised
softmax baseline (last column in \Cref{app:full}).  The results are
dataset-dependent.

On the concentric rings datasets, softmax clearly outperforms the
trained head: softmax AUROC 0.85--0.92 vs.\ head AUROC 0.51--0.67
on rings, and 0.81--0.87 vs.\ 0.50--0.66 on rings\_hard.  Similarly,
on digits (softmax 0.94--0.98 vs.\ head 0.45--0.76) and blobs
(softmax 0.72--0.79 vs.\ head 0.50--0.66), the zero-cost baseline
dominates.

However, on breast\_cancer and two\_moons, the trained head
\emph{approaches or matches} softmax discrimination.  On breast\_cancer,
\texttt{combined\_nonorm} achieves head AUROC 0.939 vs.\ softmax 0.958
--- a gap of only 0.019.  On two\_moons, \texttt{hlcc\_nonorm\_trainable}
reaches head AUROC 0.888 vs.\ softmax 0.933.  This suggests that on
tasks with well-separated classes and clear decision boundaries, the
trained head can learn meaningful per-instance discrimination, not just
the base rate.

On the harder tasks (rings, blobs), the \hlcc{}-trained head learns
the \emph{base rate} --- ``I am correct $X$\% of the time, so I should
be $X$\% confident'' --- rather than learning \emph{which specific
predictions} are likely wrong.  The softmax output, by contrast,
directly reflects the model's internal certainty about class separation
for each input, which is highly informative at small network scale.

This finding does not hold at LLM scale~\cite{mccallum2026hlcc},
where the NormShift confidence head outperforms softmax-based baselines
by +0.12 AUROC.  The likely explanation is that LLMs with 28--32
layers and 4096-dimensional hidden states provide much richer
normalisation signals, and the confidence head has access to
intermediate representations that the output logits do not capture.

\subsection{HLCC vs.\ Cross-Entropy Training Dynamics}

Training dynamics differ substantially between loss functions:
\begin{itemize}
  \item \textbf{CE converges faster.}  CE variants reach peak accuracy
        in 80--120 epochs, while HLCC variants need 140--200 epochs due
        to the indirect gradient path.
  \item \textbf{Two-phase combines benefits.}  The
        \texttt{hlcc\_norm\_2phase} variant achieves good accuracy from
        the CE warmup phase (epoch 0--50) and then develops calibrated
        confidence in the HLCC phase (epoch 50+), reaching ECE 0.077.
  \item \textbf{Confidence evolves predictably.}  In HLCC-trained
        variants, mean confidence starts at 0.5, rises steadily during
        training, and stabilises at 0.88--0.99 once accuracy plateaus.
        This matches the HLCC scoring incentive to be confident when
        correct.
\end{itemize}

% ============================================================
\section{Discussion}
\label{sec:discussion}

\paragraph{Calibration vs.\ discrimination: a dataset-dependent tension.}
A central finding of this study is that \hlcc{} produces two
distinct types of improvement whose co-occurrence depends on the
dataset.  The trained confidence head achieves excellent
\emph{calibration} across all six datasets (ECE as low as 0.020 on
two\_moons, 0.023 on breast\_cancer, 0.034 on rings).

\emph{Discrimination}, however, is dataset-dependent.  On the harder
synthetic tasks (rings, rings\_hard, blobs), the trained head achieves
worse AUROC (0.51--0.67) than the zero-cost softmax baseline
(0.72--0.92).  On breast\_cancer and two\_moons, the gap narrows
substantially: the trained head reaches AUROC 0.939 and 0.888
respectively, approaching softmax performance (0.958 and 0.933).
This suggests that the calibration--discrimination trade-off is not
inherent to \hlcc{}, but reflects the information available to the
confidence head at small network scale.

The practical implication: \textbf{for selective prediction on tasks
where the trained head approaches softmax discrimination,
\hlcc{} provides both calibration and ranking in a single pass.}
On harder tasks, softmax confidence or temperature scaling may
suffice for discrimination, while \hlcc{} adds calibration value.
At LLM scale~\cite{mccallum2026hlcc}, the richer normalisation
signals close the discrimination gap entirely.

\paragraph{OOD overconfidence is a central failure mode.}
All variants with trainable confidence heads that achieve excellent
in-distribution calibration become \emph{catastrophically}
overconfident on OOD data.  Mean confidence reaches 0.89--0.96 on
examples where accuracy is $\sim$0.29 (near random for 5 classes).
OOD AUROC drops to 0.22--0.63 --- often below random chance, meaning
the head is \emph{anti-correlated} with correctness.  By contrast,
the frozen-at-0.5 variants maintain OOD AUROC of 0.85--0.86.

This is not merely ``a risk'' but a structural consequence of the
training objective: the confidence head learns to mirror in-distribution
accuracy, and since it has never seen OOD examples during training,
it applies the same high-confidence pattern universally.  The
normalisation-derived doubt signals, which were hypothesised to
detect OOD inputs, fail to do so at small scale --- likely because
the 6 doubt signals from 2 normalisation layers lack the diversity
to distinguish distributional shift from normal variation.

This finding motivates OOD-aware training objectives, confidence
regularisation, or explicit distributional shift detectors as
necessary complements to \hlcc{} in safety-critical applications.

\paragraph{When does dataset difficulty matter?}
At 93\% accuracy (standard rings), the \hlcc{} optimal confidence
$\cstar = p/[4(1-p)] = 3.3$ saturates at 1.0, meaning the scoring
function's conservatism is never exercised.  The hard rings variant
(20\% noise, $\sim$88\% accuracy) still has high enough accuracy that
$\cstar$ saturates, but the blobs dataset ($\sim$72\% accuracy) places
predictions in the range where $\cstar \approx 0.64$, and the \hlcc{}
asymmetry between correct and incorrect predictions becomes materially
different from Brier scoring.  This suggests that \hlcc{} is most
valuable in moderate-accuracy regimes (50--80\%), which is precisely
where LLMs operate on challenging benchmarks like TruthfulQA.

\paragraph{Focal loss comparison.}
Including focal loss ($\gamma = 2$) addresses the natural question of
whether \hlcc{}'s asymmetry is simply replicating known asymmetric
loss behaviour.  The key difference: focal loss modifies the
\emph{classification} gradient (down-weighting easy examples) but
provides no mechanism for training a \emph{confidence estimator}.
\hlcc{} modifies the \emph{confidence} gradient while leaving
classification training to the base loss.  These are complementary:
the \texttt{focal\_hlcc\_norm} variant combines both, testing whether
focal's improved classification calibration compounds with \hlcc{}'s
explicit confidence training.

\paragraph{Comparison with TCP and MDCA.}
To situate \hlcc{} against established training-time calibration
methods, we compare directly with TCP~\cite{corbiere2019addressing}
and MDCA~\cite{hebbalaguppe2022stitch} using the same confidence head
architecture (hidden activations + doubt signals $\to$ 8-unit MLP
$\to$ sigmoid).  TCP replaces the \hlcc{} scoring function with
MSE$(c, \mathbb{1}[\hat{y}=y])$; MDCA replaces it with the batch-level
regulariser $|\bar{c} - \bar{a}|$.

Averaged across four datasets (5 seeds each), MDCA achieves the
lowest ECE (0.066 vs.\ \hlcc{} 0.111), confirming that directly
penalising the confidence--accuracy gap produces tighter calibration.
TCP achieves the best overall balance: ECE of 0.079, Brier 0.103,
and AUROC 0.680---outperforming \hlcc{} on all three metrics.

The result is dataset-dependent.  On two\_moons (97\% accuracy), TCP
dramatically outperforms \hlcc{} on discrimination (AUROC 0.819 vs.\
0.473; PRR 0.572 vs.\ $-$0.231), while MDCA achieves near-perfect
calibration (ECE 0.002) but poor discrimination (AUROC 0.526).  On
blobs (72\% accuracy), TCP's calibration advantage is most pronounced
(ECE 0.038 vs.\ 0.116), because the MSE target directly learns the
identity function between confidence and correctness, while \hlcc{}'s
asymmetric scoring pushes the head into a less calibrated operating
point.

This finding is important context for the paper's claims: \textbf{at
small network scale, TCP's simpler MSE objective outperforms \hlcc{}'s
asymmetric scoring function.}  The hypothesis is that \hlcc{}'s
advantage emerges at LLM scale, where (a) the confidence head must
generalise across diverse tasks rather than fit a single dataset, and
(b) the asymmetric penalty structure provides stronger gradient signal
for overconfident predictions---a regime that LLMs frequently occupy
on challenging benchmarks.

\paragraph{Gradient structure insight.}
A key architectural finding: the \hlcc{} loss cannot train classifier
weights directly, because correctness is determined by argmax (which
is non-differentiable).  Gradients from \hlcc{} flow \emph{only}
through the confidence head.  This means a CE (or focal) component is
always needed for the classifier, while \hlcc{} trains the confidence
estimator.  At LLM scale, this constraint is invisible because the
base model is frozen and only the confidence head is trained.

\paragraph{Normalisation signals at small scale.}
With only two normalisation layers and hidden dimensions of 24 and
16, the doubt signals carry enough information to improve AUROC by
0.04--0.15 in matched comparisons on the rings-family datasets.
However, the benefit is inconsistent across datasets: on
breast\_cancer, normalisation \emph{decreases} AUROC
(0.896 vs.\ 0.909), and on digits the effect is dramatic in the wrong
direction (0.506 vs.\ 0.756).  This suggests that normalisation-derived
doubt signals are most informative when the task creates varied internal
representations --- on easily separable tasks, the normalisation
statistics are too uniform to provide useful uncertainty information.
At LLM scale, where 28--32 layers provide richer normalisation dynamics,
this limitation disappears.

\paragraph{Deep ensemble comparison.}
Deep ensembles represent the gold standard for small-network
uncertainty.  The ensemble baseline (5 members, CE-trained) provides
a critical reference point: if \hlcc{} with a single network cannot
match ensembles, then the practical case for \hlcc{} at small scale
is weak.  However, \hlcc{} uses a single forward pass, whereas
ensembles require $M$ forward passes and $M\times$ the parameters,
making the comparison one of efficiency vs.\ quality.

\paragraph{Limitations.}
\begin{itemize}
  \item A four-layer feedforward network is far simpler than a
        transformer: no attention, no residual connections, no
        positional encoding.  The transferability to transformers
        rests on the LLM-scale results, not on this experiment alone.
  \item With only $\sim$1{,}000 parameters, the confidence head
        (22$\to$8$\to$1, $\sim$190 parameters) represents a
        significant fraction of total capacity.  At LLM scale, the
        confidence head is negligible relative to the base model.
  \item The OOD construction varies across datasets (geometric,
        distributional shift, Gaussian perturbation), making
        cross-dataset OOD comparisons imprecise.
  \item We do not evaluate conformal prediction, Platt scaling,
        or isotonic regression, which would provide additional
        post-hoc baseline comparisons.
  \item TCP outperforms \hlcc{} on the small-network datasets
        evaluated here.  Whether \hlcc{}'s asymmetric scoring
        provides advantages at LLM scale (where overconfidence
        on incorrect predictions is more prevalent) remains to be
        confirmed with a matched TCP comparison on large models.
\end{itemize}

% ============================================================
\section{Related Work}
\label{sec:related}

\paragraph{Calibration.}
Guo et al.~\cite{guo2017calibration} demonstrated that modern deep
networks are poorly calibrated, motivating post-hoc methods such as
temperature scaling~\cite{platt1999probabilistic} and isotonic
regression~\cite{zadrozny2002transforming}.  Naeini et
al.~\cite{naeini2015obtaining} introduced Expected Calibration Error
as a standard metric; Kumar et al.~\cite{kumar2019verified} later
showed that standard ECE is biased upward for finite samples and
proposed debiased alternatives.

\paragraph{Uncertainty estimation.}
MC Dropout~\cite{gal2016dropout} reinterprets dropout as approximate
Bayesian inference, providing uncertainty estimates at low cost.
Deep ensembles~\cite{lakshminarayanan2017simple} average predictions
across independently trained networks and remain a strong baseline.
Focal loss~\cite{lin2017focal} provides an alternative asymmetric
training signal that down-weights easy examples, improving implicit
calibration.  Mukhoti et al.~\cite{mukhoti2020calibrating} showed that
focal loss can rival or surpass temperature scaling for calibration,
and established the practice of reporting metrics before and after
post-hoc recalibration.  Zhu et al.~\cite{zhu2024calibration} provide
a comprehensive calibration benchmark across architectures and methods.

\paragraph{Auxiliary confidence networks.}
Corbi\`ere et al.~\cite{corbiere2019addressing} introduced
\emph{True Class Probability} (TCP): an auxiliary network that takes
penultimate-layer features from the main classifier and is trained
with MSE loss to predict $\mathbb{1}[\hat{y} = y]$, the probability
that the prediction is correct.  The extended version~\cite{corbiere2022confidence}
showed TCP generalises across architectures.  Our confidence head shares
TCP's architectural pattern---an auxiliary MLP on hidden
activations---but replaces the MSE target with the asymmetric \hlcc{}
scoring function and additionally feeds normalisation-derived doubt
signals.  Hebbalaguppe et al.~\cite{hebbalaguppe2022stitch}
proposed MDCA, a differentiable regulariser that penalises the
batch-level gap between average confidence and average accuracy:
$\mathcal{L}_\text{MDCA} = |\bar{c} - \bar{a}|$.  At LLM scale,
Kadavath et al.~\cite{kadavath2022language} train a P(IK)
(``probability of knowing'') head on hidden states, and
Shen et al.~\cite{shen2024thermometer} train an auxiliary MLP
to predict per-task temperature parameters.  We compare directly
against TCP and MDCA in our small-network experiments
(Section~\ref{sec:discussion}).

\paragraph{Normalisation.}
Batch normalisation~\cite{ioffe2015batch} and layer
normalisation~\cite{ba2016layer} are standard components of deep
networks.  Our work is, to our knowledge, the first to explicitly
route normalisation correction magnitudes as uncertainty signals.

\paragraph{Confidence-based assessment.}
The \hlcc{} scoring function draws on confidence-based marking in
education~\cite{mccallum2010gamedesign, gardnermedwin1995}, where
students indicate certainty alongside answers and receive
asymmetric rewards.

% ============================================================
\section{Conclusion}
\label{sec:conclusion}

We present a controlled small-network experiment isolating the
components of the \hlcc{} confidence system.  Across 13 training
configurations, 6 datasets, 4 calibration baselines, and 5 random
seeds on a 4-layer feedforward network, we find:

\begin{enumerate}
  \item \textbf{\hlcc{} loss enables calibration across all datasets.}
        Trainable confidence heads trained with CE alone show no
        calibration improvement (ECE 0.42--0.48), while
        \hlcc{}-trained variants achieve ECE as low as 0.020
        (two\_moons), 0.023 (breast\_cancer), 0.034 (rings), and
        0.068 (blobs).  This improvement is robust across 5 seeds
        and 6 datasets spanning synthetic and real-world tasks.
  \item \textbf{Discrimination is dataset-dependent.}
        On harder tasks (rings, blobs), the trained head achieves
        lower AUROC (0.51--0.67) than the softmax baseline
        (0.72--0.92).  But on breast\_cancer and two\_moons, the head
        approaches softmax: AUROC 0.939 vs.\ 0.958 and 0.888 vs.\
        0.933 respectively.  The ``softmax beats trained head''
        finding is not universal but depends on dataset structure.
  \item \textbf{Normalisation benefit is inconsistent.}  Adding
        \texttt{UncertaintyAwareLayerNorm} increases AUROC by
        0.04--0.15 on rings-family datasets but \emph{decreases}
        it on breast\_cancer and digits.  The doubt signals are most
        informative when the task creates varied internal
        representations across examples.
  \item \textbf{OOD overconfidence is a structural failure.}
        All trainable variants become catastrophically overconfident
        on OOD data (confidence 0.89--0.96, accuracy $\sim$0.29,
        often anti-correlated AUROC).  This is not a minor
        limitation but a central property of in-distribution
        confidence training that must be addressed for deployment.
  \item \textbf{Dataset difficulty matters.}  The \hlcc{} scoring
        function's conservatism only engages when accuracy is below
        $\sim$80\%.  On the blobs dataset ($\sim$72\% accuracy),
        the HLCC asymmetry is most meaningful; on easy tasks
        ($>$90\% accuracy), the optimal confidence saturates at 1.0.
  \item \textbf{The discrimination gap closes at LLM scale.}  The
        softmax-beats-trained-head finding is specific to small
        networks.  At LLM scale~\cite{mccallum2026hlcc}, the
        NormShift head outperforms logit baselines by +0.12 AUROC,
        suggesting that the value of a trained confidence head
        increases with model depth and normalisation layer count.
\end{enumerate}

The overall assessment is that \hlcc{} at small network scale is
primarily a \emph{calibration} method, with discrimination performance
that varies by dataset.  On well-separated tasks (breast\_cancer,
two\_moons), the trained head nearly matches softmax discrimination
while providing calibrated confidence.  On harder tasks (rings, blobs),
simple baselines (softmax confidence, temperature scaling, deep
ensembles) provide better per-instance ranking.  The case for \hlcc{}
strengthens with scale: the richer normalisation signals in deep
networks provide the confidence head with information that output
logits alone cannot capture.

\bibliographystyle{plainnat}
\bibliography{references}

% ============================================================
\appendix

\section{Full Results Across All Datasets}
\label{app:full}

Complete in-distribution metrics for all 13 variants across 6 datasets
(mean over 5 seeds).  Softmax AUROC provides an unsupervised baseline
for discrimination comparison.  Tables~\ref{tab:full_rings}--\ref{tab:full_blobs}
present the full results.

% --- rings ---
\begin{table}[H]
\centering
\caption{Full results for \texttt{rings} dataset (mean over 5 seeds).}
\label{tab:full_rings}
{\small
\begin{tabular}{@{}lcccccccc@{}}
\toprule
Variant & Acc & ECE$\downarrow$ & Brier$\downarrow$ & HLCC$\uparrow$ & AUROC$\uparrow$ & PRR$\uparrow$ & Conf & Softmax \\
 & & & & & & & & AUROC \\
\midrule
\texttt{ce\_nonorm}                      & .937 & .437 & .250 & 1.37 & .505 & $-$.008 & .500 & .909 \\
\texttt{ce\_norm}                        & .920 & .420 & .250 & 1.34 & .511 & .031 & .500 & .898 \\
\texttt{hlcc\_nonorm\_frozen}            & .933 & .433 & .250 & 1.37 & .513 & .021 & .500 & .919 \\
\texttt{hlcc\_norm\_frozen}              & .915 & .415 & .250 & 1.33 & .504 & .012 & .500 & .902 \\
\texttt{hlcc\_nonorm\_trainable}         & .920 & .062 & .077 & 1.66 & .584 & .244 & .937 & .879 \\
\texttt{hlcc\_norm\_trainable}           & .892 & .086 & .099 & 1.58 & .665 & .382 & .930 & .867 \\
\texttt{combined\_nonorm}                & .937 & .034 & .060 & 1.73 & .514 & .006 & .968 & .917 \\
\texttt{combined\_norm}                  & .918 & .076 & .081 & 1.68 & .600 & .205 & .970 & .909 \\
\texttt{ce\_norm\_trainable}             & .920 & .418 & .248 & 1.34 & .548 & .158 & .502 & .898 \\
\texttt{hlcc\_norm\_2phase}              & .908 & .077 & .085 & 1.63 & .639 & .294 & .949 & .886 \\
\midrule
\texttt{focal\_nonorm}                   & .921 & .421 & .250 & 1.34 & .488 & $-$.035 & .500 & .887 \\
\texttt{focal\_norm}                     & .917 & .417 & .250 & 1.33 & .516 & .051 & .500 & .855 \\
\texttt{focal\_hlcc\_norm}               & .910 & .054 & .083 & 1.61 & .603 & .189 & .904 & .852 \\
\bottomrule
\end{tabular}
}
\par\smallskip{\footnotesize\textit{Source:} \texttt{\datarepo/data/results/small\_network/results\_multi\_seed\_rings.json}}
\end{table}

% --- rings_hard ---
\begin{table}[H]
\centering
\caption{Full results for \texttt{rings\_hard} dataset (mean over 5 seeds).}
\label{tab:full_rings_hard}
{\small
\begin{tabular}{@{}lcccccccc@{}}
\toprule
Variant & Acc & ECE$\downarrow$ & Brier$\downarrow$ & HLCC$\uparrow$ & AUROC$\uparrow$ & PRR$\uparrow$ & Conf & Softmax \\
 & & & & & & & & AUROC \\
\midrule
\texttt{ce\_nonorm}                      & .912 & .412 & .250 & 1.32 & .500 & $-$.025 & .500 & .870 \\
\texttt{ce\_norm}                        & .888 & .388 & .250 & 1.28 & .544 & .117 & .500 & .839 \\
\texttt{hlcc\_nonorm\_frozen}            & .910 & .410 & .250 & 1.32 & .510 & .002 & .500 & .868 \\
\texttt{hlcc\_norm\_frozen}              & .889 & .389 & .250 & 1.28 & .518 & .040 & .500 & .836 \\
\texttt{hlcc\_nonorm\_trainable}         & .898 & .297 & .186 & 1.37 & .588 & .230 & .600 & .831 \\
\texttt{hlcc\_norm\_trainable}           & .859 & .272 & .192 & 1.30 & .662 & .362 & .588 & .813 \\
\texttt{combined\_nonorm}                & .914 & .243 & .138 & 1.45 & .565 & .126 & .671 & .865 \\
\texttt{combined\_norm}                  & .887 & .225 & .157 & 1.40 & .646 & .321 & .663 & .836 \\
\texttt{ce\_norm\_trainable}             & .888 & .386 & .248 & 1.28 & .541 & .095 & .502 & .839 \\
\texttt{hlcc\_norm\_2phase}              & .878 & .254 & .167 & 1.35 & .652 & .322 & .624 & .837 \\
\midrule
\texttt{focal\_nonorm}                   & .899 & .399 & .250 & 1.30 & .525 & .042 & .500 & .830 \\
\texttt{focal\_norm}                     & .879 & .379 & .250 & 1.26 & .528 & .077 & .500 & .820 \\
\texttt{focal\_hlcc\_norm}               & .872 & .267 & .180 & 1.33 & .657 & .339 & .605 & .817 \\
\bottomrule
\end{tabular}
}
\par\smallskip{\footnotesize\textit{Source:} \texttt{\datarepo/data/results/small\_network/results\_multi\_seed\_rings\_hard.json}}
\end{table}

% --- two_moons ---
\begin{table}[H]
\centering
\caption{Full results for \texttt{two\_moons} dataset (mean over 5 seeds).}
\label{tab:full_two_moons}
{\small
\begin{tabular}{@{}lcccccccc@{}}
\toprule
Variant & Acc & ECE$\downarrow$ & Brier$\downarrow$ & HLCC$\uparrow$ & AUROC$\uparrow$ & PRR$\uparrow$ & Conf & Softmax \\
 & & & & & & & & AUROC \\
\midrule
\texttt{ce\_nonorm}                      & .972 & .472 & .250 & 1.44 & .517 & .156 & .500 & .950 \\
\texttt{ce\_norm}                        & .970 & .470 & .250 & 1.44 & .483 & .053 & .500 & .957 \\
\texttt{hlcc\_nonorm\_frozen}            & .972 & .472 & .250 & 1.44 & .532 & .166 & .500 & .946 \\
\texttt{hlcc\_norm\_frozen}              & .970 & .470 & .250 & 1.44 & .517 & .149 & .500 & .953 \\
\texttt{hlcc\_nonorm\_trainable}         & .975 & .020 & .023 & 1.90 & .888 & .867 & .990 & .933 \\
\texttt{hlcc\_norm\_trainable}           & .968 & .030 & .032 & 1.87 & .695 & .344 & .997 & .939 \\
\texttt{combined\_nonorm}                & .972 & .028 & .028 & 1.89 & .712 & .630 & 1.000 & .946 \\
\texttt{combined\_norm}                  & .971 & .028 & .029 & 1.88 & .473 & $-$.231 & .999 & .954 \\
\texttt{ce\_norm\_trainable}             & .970 & .469 & .249 & 1.44 & .587 & .272 & .501 & .957 \\
\texttt{hlcc\_norm\_2phase}              & .966 & .032 & .034 & 1.86 & .772 & .542 & .998 & .958 \\
\midrule
\texttt{focal\_nonorm}                   & .970 & .470 & .250 & 1.44 & .527 & .167 & .500 & .953 \\
\texttt{focal\_norm}                     & .970 & .470 & .250 & 1.44 & .484 & .069 & .500 & .955 \\
\texttt{focal\_hlcc\_norm}               & .969 & .030 & .031 & 1.87 & .687 & .455 & .999 & .949 \\
\bottomrule
\end{tabular}
}
\par\smallskip{\footnotesize\textit{Source:} \texttt{\datarepo/data/results/small\_network/results\_multi\_seed\_two\_moons.json}}
\end{table}

% --- breast_cancer ---
\begin{table}[H]
\centering
\caption{Full results for \texttt{breast\_cancer} dataset (mean over 5 seeds).}
\label{tab:full_breast_cancer}
{\small
\begin{tabular}{@{}lcccccccc@{}}
\toprule
Variant & Acc & ECE$\downarrow$ & Brier$\downarrow$ & HLCC$\uparrow$ & AUROC$\uparrow$ & PRR$\uparrow$ & Conf & Softmax \\
 & & & & & & & & AUROC \\
\midrule
\texttt{ce\_nonorm}                      & .960 & .460 & .250 & 1.42 & .130 & $-$1.277 & .500 & .951 \\
\texttt{ce\_norm}                        & .958 & .458 & .250 & 1.42 & .185 & $-$1.017 & .500 & .935 \\
\texttt{hlcc\_nonorm\_frozen}            & .967 & .467 & .250 & 1.43 & .150 & $-$1.168 & .500 & .940 \\
\texttt{hlcc\_norm\_frozen}              & .963 & .463 & .250 & 1.43 & .157 & $-$1.159 & .500 & .915 \\
\texttt{hlcc\_nonorm\_trainable}         & .958 & .037 & .040 & 1.83 & .909 & .902 & .993 & .953 \\
\texttt{hlcc\_norm\_trainable}           & .951 & .047 & .048 & 1.80 & .896 & .885 & .999 & .907 \\
\texttt{combined\_nonorm}                & .963 & .023 & .033 & 1.85 & .939 & .936 & .981 & .958 \\
\texttt{combined\_norm}                  & .963 & .026 & .036 & 1.84 & .729 & .657 & .989 & .895 \\
\texttt{ce\_norm\_trainable}             & .958 & .458 & .250 & 1.42 & .440 & $-$.534 & .500 & .935 \\
\texttt{hlcc\_norm\_2phase}              & .951 & .041 & .047 & 1.80 & .889 & .874 & .992 & .890 \\
\midrule
\texttt{focal\_nonorm}                   & .965 & .465 & .250 & 1.43 & .151 & $-$1.142 & .500 & .931 \\
\texttt{focal\_norm}                     & .965 & .465 & .250 & 1.43 & .222 & $-$.949 & .500 & .911 \\
\texttt{focal\_hlcc\_norm}               & .963 & .029 & .036 & 1.83 & .823 & .794 & .979 & .850 \\
\bottomrule
\end{tabular}
}
\par\smallskip{\footnotesize\textit{Source:} \texttt{\datarepo/data/results/small\_network/results\_multi\_seed\_breast\_cancer.json}}
\end{table}

% --- digits ---
\begin{table}[H]
\centering
\caption{Full results for \texttt{digits} dataset (mean over 5 seeds).}
\label{tab:full_digits}
{\small
\begin{tabular}{@{}lcccccccc@{}}
\toprule
Variant & Acc & ECE$\downarrow$ & Brier$\downarrow$ & HLCC$\uparrow$ & AUROC$\uparrow$ & PRR$\uparrow$ & Conf & Softmax \\
 & & & & & & & & AUROC \\
\midrule
\texttt{ce\_nonorm}                      & .976 & .476 & .250 & 1.45 & .544 & .168 & .500 & .968 \\
\texttt{ce\_norm}                        & .976 & .476 & .250 & 1.45 & .483 & .002 & .500 & .976 \\
\texttt{hlcc\_nonorm\_frozen}            & .976 & .476 & .250 & 1.45 & .544 & .168 & .500 & .968 \\
\texttt{hlcc\_norm\_frozen}              & .974 & .474 & .250 & 1.45 & .452 & $-$.147 & .500 & .979 \\
\texttt{hlcc\_nonorm\_trainable}         & .973 & .050 & .031 & 1.86 & .756 & .656 & .969 & .962 \\
\texttt{hlcc\_norm\_trainable}           & .977 & .043 & .026 & 1.88 & .506 & $-$.096 & .968 & .964 \\
\texttt{combined\_nonorm}                & .976 & .048 & .028 & 1.87 & .684 & .449 & .969 & .965 \\
\texttt{combined\_norm}                  & .979 & .042 & .025 & 1.89 & .470 & $-$.186 & .968 & .975 \\
\texttt{ce\_norm\_trainable}             & .976 & .482 & .254 & 1.45 & .483 & $-$.180 & .495 & .976 \\
\texttt{hlcc\_norm\_2phase}              & .953 & .043 & .046 & 1.81 & .499 & $-$.089 & .996 & .940 \\
\midrule
\texttt{focal\_nonorm}                   & .975 & .475 & .250 & 1.45 & .564 & .252 & .500 & .953 \\
\texttt{focal\_norm}                     & .976 & .476 & .250 & 1.45 & .524 & .116 & .500 & .969 \\
\texttt{focal\_hlcc\_norm}               & .974 & .050 & .030 & 1.87 & .493 & $-$.192 & .969 & .967 \\
\bottomrule
\end{tabular}
}
\par\smallskip{\footnotesize\textit{Source:} \texttt{\datarepo/data/results/small\_network/results\_multi\_seed\_digits.json}}
\end{table}

% --- blobs ---
\begin{table}[H]
\centering
\caption{Full results for \texttt{blobs} dataset (mean over 5 seeds).}
\label{tab:full_blobs}
{\small
\begin{tabular}{@{}lcccccccc@{}}
\toprule
Variant & Acc & ECE$\downarrow$ & Brier$\downarrow$ & HLCC$\uparrow$ & AUROC$\uparrow$ & PRR$\uparrow$ & Conf & Softmax \\
 & & & & & & & & AUROC \\
\midrule
\texttt{ce\_nonorm}                      & .741 & .241 & .250 & 0.98 & .515 & .045 & .500 & .783 \\
\texttt{ce\_norm}                        & .737 & .237 & .250 & 0.97 & .514 & .051 & .500 & .755 \\
\texttt{hlcc\_nonorm\_frozen}            & .731 & .231 & .250 & 0.96 & .518 & .045 & .500 & .787 \\
\texttt{hlcc\_norm\_frozen}              & .743 & .243 & .250 & 0.99 & .518 & .063 & .500 & .756 \\
\texttt{hlcc\_nonorm\_trainable}         & .656 & .197 & .264 & 0.85 & .641 & .381 & .515 & .745 \\
\texttt{hlcc\_norm\_trainable}           & .656 & .167 & .248 & 0.85 & .654 & .348 & .531 & .735 \\
\texttt{combined\_nonorm}                & .726 & .126 & .212 & 1.03 & .662 & .404 & .710 & .784 \\
\texttt{combined\_norm}                  & .718 & .116 & .213 & 0.98 & .618 & .268 & .704 & .751 \\
\texttt{ce\_norm\_trainable}             & .737 & .234 & .249 & 0.97 & .499 & .023 & .503 & .755 \\
\texttt{hlcc\_norm\_2phase}              & .720 & .068 & .205 & 0.96 & .607 & .276 & .671 & .737 \\
\midrule
\texttt{focal\_nonorm}                   & .715 & .215 & .250 & 0.93 & .522 & .070 & .500 & .762 \\
\texttt{focal\_norm}                     & .722 & .222 & .250 & 0.94 & .510 & .049 & .500 & .739 \\
\texttt{focal\_hlcc\_norm}               & .704 & .104 & .214 & 0.94 & .618 & .284 & .660 & .720 \\
\bottomrule
\end{tabular}
}
\par\smallskip{\footnotesize\textit{Source:} \texttt{\datarepo/data/results/small\_network/results\_multi\_seed\_blobs.json}}
\end{table}

\section{Baseline Comparison}
\label{app:baselines}

Temperature scaling, MC dropout (30 forward passes), and deep ensemble
(5 members) baselines were evaluated on the breast cancer dataset.
Focal loss ($\gamma = 2$) variants are included as rows 11--13 in all
result tables above.  Full baseline comparison data is available at:

\par\smallskip{\footnotesize\textit{Source:} \texttt{\datarepo/data/results/small\_network/baselines\_breast\_cancer.json}}

\section{Hyperparameter Sensitivity}
\label{app:hyperparam}

Sensitivity analysis for learning rate, HLCC alpha parameter, noise
rate, focal gamma, and early stopping patience will be included here.

\end{document}
